% Figure 1: glycoMusubi resource and hierarchical benchmark overview
% Publication-quality TikZ — Bioinformatics / Nature style
\begin{figure*}[!ht]
\centering

% ═══════════════════════════════════════════════════════════════════
% Semantic colour palette (restrained: 3 hues + neutrals)
% ═══════════════════════════════════════════════════════════════════
\definecolor{cPrimary}{HTML}{2b6a99}     % muted steel blue — infrastructure
\definecolor{cPrimaryLight}{HTML}{d6e6f0} % very pale blue tint
\definecolor{cAccent}{HTML}{c04e3f}       % muted warm red — boundary / alert
\definecolor{cAccentLight}{HTML}{f5dcd8}  % pale red tint
\definecolor{cDark}{HTML}{1d1d2c}         % near-black for text & outlines
\definecolor{cMid}{HTML}{5a6068}          % medium gray for secondary text
\definecolor{cLight}{HTML}{b0b5bb}        % light gray for hairlines
\definecolor{cFaint}{HTML}{e8eaed}        % very faint gray for grid
\definecolor{cWhite}{HTML}{ffffff}
% KG node tones — two-tone scheme: primary entities vs. secondary
\definecolor{cNodeA}{HTML}{2b6a99}        % primary entity (protein, enzyme, glycan)
\definecolor{cNodeB}{HTML}{6b97b5}        % secondary entity

% ═══════════════════════════════════════════════════════════════════
% Global style definitions
% ═══════════════════════════════════════════════════════════════════
\begin{tikzpicture}[
    x=1cm, y=1cm,
    % ── reusable styles ──
    panellabel/.style={font=\fontsize{10pt}{12pt}\selectfont\bfseries,
                       text=cDark, anchor=north west},
    seclabel/.style={font=\fontsize{6.5pt}{8pt}\selectfont,
                     text=cMid, anchor=south},
    srcbox/.style={draw=cLight, rounded corners=2pt, fill=cPrimaryLight,
                   minimum width=2.0cm, minimum height=0.48cm,
                   font=\fontsize{6pt}{7pt}\selectfont\bfseries,
                   text=cPrimary, line width=0.35pt, inner sep=0pt},
    etlbox/.style={draw=cPrimary!70!black, rounded corners=2pt,
                   fill=cPrimary, text=white,
                   minimum width=1.6cm, minimum height=0.48cm,
                   font=\fontsize{6pt}{7pt}\selectfont\bfseries,
                   line width=0.35pt, inner sep=0pt},
    kgnode/.style={circle, minimum size=0.78cm,
                   font=\fontsize{4pt}{5pt}\selectfont\bfseries,
                   text=white, inner sep=0pt,
                   line width=0pt},
    thinedge/.style={cLight, line width=0.3pt},
    connector/.style={-{Stealth[length=1.4mm, width=0.9mm]},
                      cLight, line width=0.4pt},
    thickconn/.style={-{Stealth[length=1.6mm, width=1mm]},
                      cMid, line width=0.55pt},
    % bar chart
    barL1/.style={fill=cPrimary},
    barL2/.style={fill=cAccent},
    ticklabel/.style={font=\fontsize{5.5pt}{7pt}\selectfont, text=cMid},
    metriclabel/.style={font=\fontsize{5.5pt}{7pt}\selectfont, text=cDark},
    tasklabel/.style={font=\fontsize{6pt}{7pt}\selectfont, text=cDark},
    subtasklabel/.style={font=\fontsize{5pt}{6pt}\selectfont, text=cLight},
]

% ─────────────────────────────────────────────────────────────────
% PANEL (a) — top half: resource construction
% ─────────────────────────────────────────────────────────────────
\node[panellabel] at (-0.2, 7.0) {\textsf{a}};

% ── Section labels ──
\node[seclabel] at (1.0, 6.65) {Data sources};
\node[seclabel] at (3.7, 6.65) {ETL pipeline};
\node[seclabel] at (8.5, 6.65) {Knowledge graph};

% ── Thin horizontal rule beneath section labels ──
\draw[cFaint, line width=0.3pt] (-0.1, 6.55) -- (11.3, 6.55);

% ── Source column (left) ──────────────────────────────────────
\foreach \i/\name in {0/GlyGen, 1/GlyTouCan, 2/UniProt,
                       3/ChEMBL, 4/PTMCode, 5/Reactome}{
    \node[srcbox] (src\i) at (1.0, {6.1 - \i*0.62}) {\name};
}

% ── Merge connector: sources → bus → ETL ──────────────────────
% Vertical collector line
\draw[cLight, line width=0.6pt] (2.55, 2.95) -- (2.55, 6.1);
% Horizontal stubs from each source
\foreach \i in {0,...,5}{
    \draw[cLight, line width=0.3pt] (src\i.east) -- (2.55, {6.1 - \i*0.62});
}
% Trunk arrow → ETL
\draw[thickconn] (2.55, 4.52) -- (2.75, 4.52);

% ── ETL pipeline (centre) ────────────────────────────────────
\foreach \j/\stage in {0/Download, 1/Clean, 2/Build, 3/Validate}{
    \node[etlbox] (etl\j) at (3.7, {5.85 - \j*0.70}) {\stage};
}
% Stage arrows
\foreach \j in {0,1,2}{
    \pgfmathtruncatemacro{\jnext}{\j+1}
    \draw[connector, cPrimary!60] (etl\j) -- (etl\jnext);
}
% Enclosing hairline box
\draw[cLight, rounded corners=3pt, line width=0.3pt, dashed]
    (2.72, 3.28) rectangle (4.68, 6.20);

% ── Output arrow: ETL → KG ───────────────────────────────────
\draw[thickconn] (4.68, 4.74) -- (5.55, 4.74);

% ── Knowledge graph (right, focal element) ────────────────────
% 10 nodes on an ellipse — centre (8.5, 4.2), rx=2.35, ry=2.05
\pgfmathsetmacro{\cx}{8.5}
\pgfmathsetmacro{\cy}{4.30}
\pgfmathsetmacro{\rx}{2.35}
\pgfmathsetmacro{\ry}{2.05}

% Pre-draw edges (behind nodes)
% Define node angles: glycan=90, protein=54, enzyme=18, site=-18,
%   disease=-54, compound=-90, motif=-126, reaction=-162,
%   variant=162, cell_loc=126
\foreach \angA/\angB in {90/54, 18/90, 54/-54, 54/-18, 54/162,
                          18/-90, 90/-126, 18/-162, 54/126}{
    \draw[cFaint, line width=0.25pt]
        ({\cx+\rx*cos(\angA)}, {\cy+\ry*sin(\angA)}) --
        ({\cx+\rx*cos(\angB)}, {\cy+\ry*sin(\angB)});
}

% Draw nodes
\foreach \ang/\col/\lab in {
    90/cNodeA/glycan,
    54/cNodeA/protein,
    18/cNodeA/enzyme,
    -18/cNodeB/site,
    -54/cNodeB/disease,
    -90/cNodeB/compound,
    -126/cNodeB/motif,
    -162/cNodeB/reaction,
    162/cNodeB/variant,
    126/cNodeB/cell loc.%
}{
    \node[kgnode, fill=\col] at ({\cx+\rx*cos(\ang)}, {\cy+\ry*sin(\ang)}) {\lab};
}

% ── Scale summary (elegant single line beneath graph) ─────────
\node[font=\fontsize{5.5pt}{7pt}\selectfont, text=cMid, anchor=north]
    at (\cx, {\cy - \ry - 0.40})
    {78,263 nodes\enspace\textbullet\enspace 2.5\,M edges\enspace\textbullet\enspace
     10 types\,/\,14 relations};

% ── Thin separator between panels ─────────────────────────────
\draw[cFaint, line width=0.25pt] (-0.1, 1.65) -- (11.3, 1.65);

% ─────────────────────────────────────────────────────────────────
% PANEL (b) — bottom half: hierarchical benchmark
% ─────────────────────────────────────────────────────────────────
\node[panellabel] at (-0.2, 1.40) {\textsf{b}};

% Layout: bars span x = [2.0, 9.0] → 7cm = performance 0–1
\pgfmathsetmacro{\bL}{2.0}   % bar left edge
\pgfmathsetmacro{\bW}{7.0}   % bar total width (=1.0)
\pgfmathsetmacro{\bH}{0.32}  % bar height
\pgfmathsetmacro{\bG}{0.60}  % bar gap (centre-to-centre spacing)

% ── Faint grid ──
\foreach \g in {0, 0.2, 0.4, 0.6, 0.8, 1.0}{
    \draw[cFaint, line width=0.2pt]
        ({\bL + \g*\bW}, -3.55) -- ({\bL + \g*\bW}, 1.15);
}

% ── Bars ──
% Macro: \drawbar{index}{task}{subtitle}{metric}{value}{style}
% y = 0.85 - index * bG
\newcommand{\drawbar}[6]{%
    \pgfmathsetmacro{\yc}{0.85 - #1*\bG}%
    \pgfmathsetmacro{\yt}{\yc + \bH/2}%
    \pgfmathsetmacro{\yb}{\yc - \bH/2}%
    \pgfmathsetmacro{\xr}{\bL + #5*\bW}%
    \fill[#6] (\bL, \yb) rectangle (\xr, \yt);%
    % task label (left)
    \node[tasklabel, anchor=east] at ({\bL - 0.12}, \yc) {#2};%
    % subtitle
    \node[subtasklabel, anchor=east] at ({\bL - 0.12}, {\yc - 0.17}) {#3};%
    % metric (right of bar)
    \node[metriclabel, anchor=west] at ({\xr + 0.10}, \yc) {#4};%
}

\drawbar{0}{N-linked classification}{Binary}{F1 = 0.924}{0.924}{barL1}
\drawbar{1}{Site ranking}{Per protein}{R@3 = 0.683}{0.683}{barL1}
\drawbar{2}{Structural family}{8 classes}{Top-2 = 0.858}{0.858}{barL1}
\drawbar{3}{Structural cluster}{20 clusters}{H@5 = 0.786}{0.786}{barL1}

% ── Determination boundary ──
\pgfmathsetmacro{\bndY}{0.85 - 3.5*\bG}
% Accent-tinted band
\fill[cAccentLight, opacity=0.35]
    (-0.1, {\bndY - 0.12}) rectangle (11.3, {\bndY + 0.12});
% Hairline
\draw[cAccent, line width=0.5pt] (\bL - 0.3, \bndY) -- ({\bL + \bW + 0.3}, \bndY);
% Label
\node[font=\fontsize{5pt}{6pt}\selectfont\itshape, text=cAccent,
      anchor=east]
    at ({\bL + \bW + 0.3}, {\bndY + 0.18}) {determination boundary};

% ── Below-boundary bar (accent colour) ──
\drawbar{4}{Exact glycan ID}{2{,}357 glycans}{MRR = 0.066}{0.066}{barL2}

% ── X-axis ──
\pgfmathsetmacro{\axY}{0.85 - 4*\bG - \bH/2 - 0.15}
\draw[cMid, line width=0.35pt] (\bL, \axY) -- ({\bL+\bW}, \axY);
\foreach \t/\l in {0/0, 0.2/0.2, 0.4/0.4, 0.6/0.6, 0.8/0.8, 1.0/1.0}{
    \draw[cMid, line width=0.25pt]
        ({\bL+\t*\bW}, \axY) -- ({\bL+\t*\bW}, {\axY - 0.06});
    \node[ticklabel, anchor=north] at ({\bL+\t*\bW}, {\axY - 0.08}) {\l};
}
\node[font=\fontsize{6pt}{7pt}\selectfont, text=cMid, anchor=north]
    at ({\bL + \bW/2}, {\axY - 0.38}) {Performance};

% ── Specificity arrow (left margin, spanning all bars) ──
\pgfmathsetmacro{\specTop}{0.85 + \bH/2}
\pgfmathsetmacro{\specBot}{0.85 - 4*\bG - \bH/2}
\draw[-{Stealth[length=1.2mm, width=0.8mm]}, cLight, line width=0.35pt]
    (0.35, \specTop) -- (0.35, \specBot);
\node[font=\fontsize{4.5pt}{5.5pt}\selectfont, text=cLight,
      rotate=90, anchor=south] at (0.20, {(\specTop+\specBot)/2})
    {increasing specificity};

\end{tikzpicture}

\caption{\textbf{glycoMusubi resource and hierarchical benchmark overview.}
\textbf{(a)}~Data integration pipeline: six public databases are processed
through an automated four-stage ETL pipeline to construct a heterogeneous
knowledge graph with 78{,}263 nodes (10~types) and 2{,}508{,}960 edges
(14~relation types).
\textbf{(b)}~Hierarchical prediction benchmark overview: performance across five
tasks of increasing specificity.  The determination boundary marks the sharp
performance drop from structural-class prediction (H@5\,=\,0.786) to exact
glycan identification (MRR\,=\,0.066). Fig.~\ref{fig:twolayer}c expands this
view with all seven tasks annotated by Layer~1/2 assignment.}
\label{fig:kg_benchmark}
\end{figure*}
